\section{Introduction}
Instance image retrieval ranks gallery images depicting the same particular object or landmark as a query. A global descriptor enables inexpensive similarity search, but its usefulness depends on both the learned representation and the processing applied after the initial ranking. Fine-tuning can adapt a classification backbone to retrieval, while query expansion, local verification, fusion, and graph propagation offer additional ways to alter rankings without updating encoder weights.

The VGG16--GeM pipeline provides an established setting for examining these choices. VGG networks supply convolutional representations~\cite{vgg}, and generalized-mean pooling (GeM) supports compact descriptors learned from structure-from-motion (SfM) supervision~\cite{gem}. However, adding augmentation, attention, or more epochs does not by itself establish an improvement. Likewise, a search method that helps a fixed final checkpoint need not outperform a stronger historical checkpoint. Mixing these comparisons can obscure what an experiment actually demonstrates.

We examine three available training histories: a baseline reproduction, denoted B0, an augmentation-only run (\Rone), and an augmentation-plus-spatial-gate run (\Rtwo). The training comparison is historical rather than a fully controlled ablation. We separately conduct a fixed-encoder search study using the two new epoch-100 checkpoints, re-extracted with the same FP32 inference protocol. A third evaluation applies five fixed methods to validation-selected B0 on both clean and corrupted queries, so search conclusions also include the stronger baseline checkpoint.

Our study addresses three questions. \emph{RQ1:} How do validation-selected and final checkpoints compare, and what evidence supports the training modifications? \emph{RQ2:} Which post-training methods improve clean and corrupted-query retrieval for the same fixed encoder? \emph{RQ3:} How much query-time computation is required for these gains?

The contributions are empirical: (1) an audit of training configuration, data splits, and checkpoint selection; (2) a four-family search study for two final encoders, extended by a five-method clean/robustness evaluation of validation-selected B0; and (3) a comparison of retrieval-only and direct end-to-end latency, retaining unsuccessful methods and exploratory sensitivity results. The work evaluates combinations of established components; it does not introduce a new pooling, attention, or geometric-verification algorithm.
